\section{讨论与结论
\label{sec:Conclusions}
}

本文给出了 CT18 族部分子分布函数（PDF），包括 CT18Z、CT18A 与 CT18X 替代拟合。CT18 是 CTEQ-TEA 组全球分析给出的下一代 NNLO（以及 NLO）
质子 PDF。它承接 CT14 与 \CTHERAII~NNLO 分布发布之后的下一次更新，后者正是 CT14 发表后因 HERA I 与 II 精密合并数据发布而催生。
CT18 是名义上的 CTEQ-TEA PDF 组合，我们建议一切一般用途采用之。CT18A 是把 ATLAS 7 TeV $W/Z$ 数据~\cite{Aaboud:2016btc} 加入 CT18 拟合的产物；CT18X 是 CT18 的一个变体，对低 $x$ DIS 数据采用依赖 $x$ 的 QCD 标度（连同稍大的 1.4 GeV 粲夸克质量）；
CT18Z 则囊括上述全部改动，一般而言与 CT18 差异最大。CT18 与 CT18Z 数据集之间的差异仅在少数应用场合不可忽略；在这些场合需把差异折入总 PDF 不确定度，例如取 CT18 与 CT18Z 不确定度的包络。CT18A 可用于更全面地考察奇异夸克分布的不确定度范围。类似地，低 $x$ 重求和的可能影响可用 CT18X 探索。

尽管部分早期 7、8 TeV LHC Run-1 数据——包括矢量玻色子
\cite{Aaij:2012vn, Chatrchyan:2013mza, Chatrchyan:2012xt, Aad:2011dm} 与喷注 \cite{Aad:2011fc, Chatrchyan:2012bja} 的单举产生测量——曾被用作 CT14 拟合的输入，CT18 才是第一个实质性纳入 LHC 完整 Run-1 最重要实验数据的 CT 分析，涵盖 7、8 TeV 矢量玻色子、喷注与顶夸克对的单举产生测量。
CT18 全球分析所含具体数据集的详细信息见表s.~\ref{tab:EXP_1} 与~\ref{tab:EXP_2}，新纳入的数据见后表。
随着 LHC 测量精度迅速提升，全球分析的重心已转向利用最先进的理论计算，在 LHC 数据覆盖的宽广 $x$、$Q$ 范围内提供精确预言。实现这一目标需要在理论、实验与统计领域的长期多管齐下的努力。

来自实验信息一侧的挑战，是全球分析中对相关且相容数据集的挑选与实现。具体而言，我们纳入了对感兴趣的 PDF 具有敏感度、且有 NNLO 预言可得的过程。
例如对 ATLAS 喷注数据，我们利用 ATLAS 去关联模型尽量纳入尽可能大的快度区间，而非只用单一快度区间。我们注意到只用单一
快度区间可能带来选择偏差。由于 ATLAS 喷注数据中残留的张力，其结果可能是较大的 $\chi^2 / N_\mathit{pt}$，且相对 CMS 喷注数据 PDF 敏感度降低，参见 Sec.~\ref{sec:DataJets}。
类似地，为把 $t \bar t$ 微分截面测量并入 CT18 全球分析，
我们采用两个 ATLAS $t \bar t$ 单微分可观测量（利用统计关联）与 CMS 的双微分测量，以尽可能纳入更多信息。同样，若只用单一微分分布会有偏差风险；然而一些 $t\bar t$ 可观测量彼此存在张力，参见 Sec.~\ref{sec:DataTop}。
CT18 全球分析表明，CT14 全球分析中的既有数据集仍具有非常强的拉动，倾向于稀释新数据的影响。
例如，低能 DIS 与 Drell-Yan 数据、精密 HERA 数据以及 D\O{} 在 $9.7\mbox{ fb}^{-1}$
处的电子电荷不对称精确测量~\cite{D0:2014kma} 对于探查仅凭 LHC Run 1 数据无法分辨的夸克味组合仍然重要。此外，大多数实验测量含有
大量关联系统不确定度；我们在考察这些测量的 PDF 影响时已计及这些系统误差。
另外，我们还考察了重要 LHC 过程的 PDF 误差，并检验了 Hessian 与 LM 途径的一致性。

来自理论方面的挑战有三：审视理论预言对 QCD 标度选择的依赖并与实验精度对比；探索所选非微扰 PDF 参数化形式对全球分析及不确定度的影响；以及能够进行快速而准确的理论计算。
在名义 CT18 拟合中，我们采用规范的
QCD 重正化与因子化标度选择，这通常能稳定高阶理论修正。在研究 PDF 不确定度时考虑了备选标度选择的拟合，参见 Sec.~\ref{sec:FinalPDFuncertainty}。

特别是对 CT18 NNLO PDF，我们对精密 DIS、Drell-Yan、喷注与 $t \bar t$ 过程一贯采用 NNLO 计算，参看 Sec.~\ref{sec:Theory}。各过程具体的 QCD 标度选择列于表s.~\ref{Theory-Calc-II} 与~\ref{Theory-Calc-VB}。
例如，在低喷注横动量处，分别取单举喷注或
领头喷注横向动量为动量标度的 NNLO 理论预言之间存在不可忽略的差异~\cite{Currie:2018xkj}。CTEQ-TEA 组采用的名义选择是单举喷注 $p_T$。
我们观察到，即便在这两种标度选择预言差异重要的运动学区域，拟合出的胶子 PDF 对此也不太敏感。
全球拟合的这种稳健性，一方面因为其他数据也在相应运动学区域约束胶子 PDF，另一方面可能因为可观的系统不确定度带来的补偿效应。
为与 LHC 高精度数据比较，理论预言还必须计入电弱修正。细节见 Sec.~\ref{sec:calcs}，并参考表~\ref{tab:EWcorrections}。


为考察拟合对演化起始标度 $Q_0$（约 1.3 GeV）处 PDF 非微扰函数形式的依赖，我们抽样了一大组——$\mathcal{O}(250)$ 个——候选拟合形式，它们的拟合参数数目相当。（采用更灵活的参数化以更好捕捉 PDF 的 $x$ 依赖性变化，参看 Appendix \ref{sec:AppendixParam}。）
研究结果见图~\ref{fig:params}。
随着全球分析中拟合参数数目增加，
我们通常观察到 $\chi^2$ 的稳步改善；只要被拟合参数 $\lesssim\! 30$ 个，改善一般持续增长，超过此数后扩展参数化试图描述统计噪声，拟合便趋于不稳定。


为执行与{\it 单个数据}统计误差可小至 0.1\% 的精密数据比较的 NNLO CT18 全球拟合，
我们需要高数值精度的快速理论计算。
因此，在结构函数与截面计算中使用各种快速接口成为必需与惯例。
为此，我们在内部开发了 QCD 相互作用 NNLO 精度的快速 ApplGrid 与 FastNLO 计算，参看 Sec.~\ref{sec:Theory}。
此外，我们还并行化了全球拟合算法以大大加速收敛，讨论见 Appendix~\ref{sec:AppendixCodeDevelopment}。


LHC 各实验合作组成功采集了大量高精度数据。检验与这些精密数据的符合需要统计方法学的进步。
哪些合格的 LHC 实验能为 CTEQ-TEA PDF 提供有前景的约束？
LHC 实验彼此之间以及与其他实验是否一致？
四种方法的强大组合给出了一致的答案：
1) \texttt{PDFSense} 与 $L_2$ 敏感度；
2) \texttt{ePump} 程序；
3) 有效高斯变量；
4) LM 扫描。
后两种方法在此前的 CTEQ-TEA 全球分析（如 CTEQ6、CT10 与 CT14）中已经引入，前两项技术则是在 CT18 全球分析过程中发明的。

\texttt{PDFSense} 程序~\cite{Wang:2018heo} 提供了一种简便方式，可在 $x$-$Q$ 平面上可视化数据对 PDF 的潜在影响。此外，简单的 $L_2$ 敏感度变量~\cite{Hobbs:2019gob} 有助于探索不同实验之间的符合程度，功能类似 LM 扫描，但在全 $x$ 或 $Q$ 范围上使用快得多的 Hessian 形式体系。示例见 Secs.~\ref{sec:L2} 与 \ref{sec:AppendixCT18Z}。

互补的 \texttt{ePump} 程序~\cite{Schmidt:2018hvu} 含有估计新数据对最优拟合与 Hessian 误差 PDF 影响的快速高效方法。针对此前 CT14 全球拟合的大量验证亦已完成~\cite{Hou:2019gfw}。
上述四种技术在 CT18 分析中的应用贯穿全文。

在这四个 CT18 拟合中，我们发现对 PDF 有重要影响的测量包括：ATLAS 与 CMS 单举喷注产生、LHCb $W,Z$ 矢量玻色子产生，以及 ATLAS $\sqrt{s}=8$ TeV $Z$ 玻色子横向动量数据。我们发现半单举（SI）DIS 实验（如 NuTeV 与 CCFR 双缪子产生）一方与某些 LHC 实验——特别是 ATLAS 7 TeV $W/Z$ 产生测量、在一定程度上还有 LHCb $W/Z$ 测量——另一方，对奇异夸克 PDF 存在相互矛盾的偏好。对 LHC 测量与理论预言进行基准比对，加上新的 (SI)DIS 实验，可以非常有效地化解这些张力。展望未来，为配合 HL-LHC 的发现计划，要达到这些规划实验的终极精度，还需要持续努力去驾驭对撞机数据中的实验张力。我们预见理论方法与数据分析方法（包括本研究用于探究数据相容性的方法）之间的互动，以及额外的高精度实验（如电子-离子对撞机（EIC）之类的高亮度 DIS 对撞机~\cite{Accardi:2012qut}），对取得这样的进步不可或缺。

与 CT14 相比，新数据与理论进展使 CT18 发生了如下变化：
1) $x \sim 0.3$ 处 $g(x,Q)$ 变小（主要缘于 ATLAS 与 CMS 7 TeV 喷注数据及 ATLAS 8 TeV $Z$ $p_T$ 数据；CMS 7 与 8 TeV 喷注数据之间发现一定张力）；
2) 小 $x$ 处 $u$、$d$、$\bar u$、$\bar d$ 有所变化，例如 $x \sim 0.2$ 处更大的 $d$ 与 $d/u$、更小的 $\bar d / \bar u$（主要缘于 LHCb $W,Z$ 快度数据与 CMS 8 TeV $W$ 轻子电荷不对称数据）；
3) 小 $x$ 处更大的 $s$ 与 $(s+\bar s)/(\bar u + \bar d)$（主要缘于 LHCb $W,Z$ 快度数据，并在 CT18A 与 CT18Z 拟合中被 ATLAS 7 TeV $W/Z$ 数据进一步增强）。
虽然单个 $t \bar t$ 数据点的敏感度可与 LHC 单个喷注数据点相当，但由于 $t \bar t$ 数据点数目少，其总敏感度较小。因此我们没有发现纳入 CT18 分析的 $t \bar t$ 双微分分布带来可察觉的影响。文献~\cite{Hou:2019jxd} 在一项 CT14 分析中也报告了类似发现。

尽管中心预言有这些变化，CT18 NNLO PDF 在各自误差带内仍与 CT14 NNLO 一致。
CT18 与 CT14 PDF 比较以及对数据拟合质量的更多细节见 Secs.~\ref{sec:OverviewCT18} 与 \ref{sec:Quality}。

CT18 预言对唯象可观测量的若干影响已在 Sec.~\ref{sec:StandardCandles} 综述。
与 CT14 NNLO 的计算相比，CT18 中 $gg \rightarrow H$ 与 $t\bar t$ 总
截面均略有下降。$W$ 与 $Z$
截面虽仍与 CT14 一致，但因奇异性增强而略有上升。Sec.~\ref{sec:PDFratios} 所示 CT14 NNLO 的常见奇异/非奇异 PDF 比值，在 PDF 不确定度内与 ATLAS 和 CMS 的独立测定一致。

我们还给出了 CT18 全球拟合对 $\alpha_s$ 数值的含义，见 Sec.~\ref{sec:AlphasDependence}。
完整 CT18 数据集在 NNLO 偏好 $\alpha_s(M_Z)\! =\! 0.1164\! \pm\! 0.0026$（68\% C.L.）。CT18Z 的相应值基本相同：$0.1169\pm0.0027$。
相比之下，纳入很少 LHC 数据的 CT14 测定为 90\% C.L. 下 $\alpha_s(M_Z)\! =\! 0.115^{+0.006}_{-0.004}$。

对粲夸克（极点）质量 $m_c$ 的 LM 扫描（图s.~\ref{fig:lm_mc} 与~\ref{fig:lm_mcZ}）支持在 CT18 与 CT18Z 拟合中分别使用 1.3 GeV 与 1.4 GeV。
值得注意的是，HERA 合并粲数据偏好略小的 $m_c$ 值，而 CT18Z 拟合中的 ATLAS 7 TeV $W/Z$ 数据偏好更大的 $m_c$ 值。关于被拟合粲贡献对 LHC $W/Z$ 截面预言影响的评论见 Sec.~\ref{sec:ellipse} 末尾。
与其他组 PDF 相关的部分子光度及预言的比较见 Secs.~\ref{sec:PDFLuminosities}、\ref{sec:moments} 与 \ref{sec:ATL7ZWchi2}。

为便于与格点 QCD 社区的结果直接比较，我们还在 Sec.~\ref{sec:moments} 给出 CT18 对各种 PDF 矩与求和规则的预言。
总体而言，对大多数格点可观测量，CT18 与其他唯象拟合工作的结果相符良好。目前系统性效应使得许多格点计算明显超出当代唯象学的预言，唯一的例外是胶子矩 $\langle x \rangle_g$——格点结果相对 PDF 拟合偏低。我们期待格点模拟与 PDF 唯象学的互补进展在未来几年改善这一状况，并为确定核子纵向结构的协同 PDF-格点事业~\cite{Hobbs:2019gob,Lin:2017snn} 铺平道路。

最终的 CT18 PDF 以 1 套中心加 58 套
Hessian 本征矢组的形式在 NNLO 与 NLO 发布。物理可观测量的 90\% C.L. PDF
不确定度可用对称~\cite{Pumplin:2002vw} 或非对称
~\cite{Lai:2010vv,Nadolsky:2001yg} 主公式、把本征矢对的贡献按平方和相加来估算。这些 PDF
针对中心 QCD 耦合 $\alpha_s(M_Z)=0.118$ 确定，
与世界平均 $\alpha_s$ 值一致。为估计
组合的 PDF+$\alpha_s$ 不确定度，我们对 $\alpha_s(M_Z)=0.116$ 与 0.120 另提供两套最优拟合组。$\alpha_s(M_Z)$ 引起的
90\% C.L. 变动可取两套 $\alpha_s$ 组预言之差的二分之一来估计。90\% C.L.、含关联的
PDF+$\alpha_s$ 不确定度也可由 PDF 不确定度与 $\alpha_s$
不确定度按平方和叠加得到~\cite{Lai:2010nw}。
除了这些通用 PDF 组之外，我们还提供 $\alpha_s(M_Z)=0.111-0.123$ 的一系列 (N)NLO
组合，以及采用非标准 5 味方法的重夸克方案（活跃味数分别为 3、4、6）的附加组合。

CT18 PDF 组的参数化文件以独立形式通过 CTEQ-TEA 网站~\cite{CT18website} 分发，或作为
LHAPDF6 库~\cite{LHAPDF6} 的一部分提供。为向后兼容 LHAPDF
5.9.X 版本，我们的网站也提供 LHAPDF5 格式的 CT18 网格，以及 LHAPDF5 库 CTEQ-TEA 模块的更新版——编译时必须一并包含，以支持调用 CT18 附带的全部本征矢组~\cite{LHAPDF5}。





\section*{致谢}
我们深深感谢我们的朋友与同事 Jon Pumplin 数十年富有成效的合作，以及他在本文中所作的、对 CTEQ-TEA 全球分析的最后一项重大贡献。
感谢 Stefano Camarda、Amanda Cooper-Sarkar、Alexander Glazov、Lucian Harland-Lang、Jan Kretzschmar、Bogdan Malescu、Wally Melnitchouk、Dave Soper 以及各位 CTEQ 同仁富有启发的讨论。本工作得到美国能源部 Grants No.~DE-SC0010129（SMU）与 No.~DE-FG02-95ER40896（匹兹堡大学）的部分支持；美国国家科学基金会
Grants No.~PHY-1719914 与 No.~PHY-2013791（MSU）及 No.~PHY-1820760（匹兹堡大学）的支持，以及 PITT PACC 的部分支持。T.~J.~Hobbs 感谢 JLab EIC Center Fellowship 的支持。
M.Guzzi 的工作由国家科学基金会 Grant No.~PHY-1820818 支持。J.G. 的工作受国家自然科学基金（NNSF）
Grants No.~11875189 与 No.~11835005 支持，S.D. 与 I.S. 的工作受 NNSF Grants No.~11965020 与 No.~11847160 支持。C.-P.~Yuan 还感谢粒子物理 Wu-Ki Tung 冠名讲席教授职位的支持。
